\chapter{Suivi de tendance}\label{ch:s1:trend-following}

Prenez tous les marchés à terme liquides, prenez une position longue sur ceux qui ont monté au cours de l'année écoulée et une position courte sur ceux qui ont baissé, et dimensionnez chaque position de manière que chacune
porte le même risque. Hurst, Ooi et Pedersen ont appliqué cette règle aux marchés mondiaux depuis 1880 et ont constaté des rendements moyens positifs dans chaque décennie, une faible corrélation
avec les actions et les obligations, et de bonnes performances dans huit des dix pires drawdowns d'un portefeuille 60/40. C'est la stratégie systématique la plus simple qui soit,
et l'une des rares à disposer d'un siècle d'observations. Sur l'univers synthétique de quarante contrats à terme de ce chapitre, la règle de douze mois rapporte un ratio de Sharpe
de 0.66 après coûts sur trente ans, et, lors des deux grands krachs boursiers, le portefeuille de \omterm{def:s1:trend-following:trend}{suivi de tendance} gagne 21~\% et 39~\% tandis que les actions perdent 24~\% et 35~\%. La construction
repose sur deux composantes~: \texttt{firm.synthfut}, l'univers synthétique de contrats à terme sur lequel ce chapitre et les six suivants s'appuient, et \texttt{firm.trendfollow}.

\begin{definition}[Suivi de tendance, momentum temporel]\label{def:s1:trend-following:trend}
Le \emph{suivi de tendance}\index{suivi de tendance} (trend following) consiste à prendre position sur un marché dans le sens de sa variation récente, à l'achat après les hausses et à la vente après
les baisses, et ce sur de nombreux marchés à la fois. Le \emph{momentum temporel}\index{momentum temporel} (time-series momentum) est sa règle la plus simple~: la position sur chaque marché a
le signe du rendement en excès passé de ce marché sur une fenêtre de $L$ jours, quoi qu'aient fait les autres marchés~; le \omterm{def:s1:momentum:cross}{momentum transversal} (chapitre~5) classe au contraire
les marchés les uns par rapport aux autres.
\end{definition}

\section{Familles de signaux~: rendements, moyennes mobiles, cassures}

Notons $r_{it}$ le rendement quotidien en excès du marché $i$, c'est-à-dire le rendement d'une position reportée sur le contrat à terme le plus proche, et $P_{it} =
\sum_{s \le t} r_{is}$ sa somme cumulée, un logarithme du prix qui ignore les écarts de report (série continue de contrats à terme du livre~7, chapitre~2). Trois familles de
règles transforment $P$ en une direction, chacune avec une vitesse propre.
\begin{itemize}
\item \textbf{Rendement passé.} $s_{it} = \operatorname{sign}(P_{it} - P_{i,t-L})$~: \omterm{def:s1:trend-following:trend}{momentum temporel} de fenêtre $L$.
\item \textbf{Croisement de moyennes mobiles.} $s_{it} = \operatorname{sign}(\bar P^{(n)}_{it} - \bar P^{(m)}_{it})$, la moyenne courte de $P$ sur $n$ jours
comparée à la moyenne longue sur $m > n$ (livre~7, chapitre~7).
\item \textbf{Cassure.} La position passe à l'achat lorsque $P$ atteint son plus haut niveau des $L$ derniers jours, à la vente lorsqu'il atteint son plus bas, et elle est maintenue sinon.
\end{itemize}

\begin{definition}[Règle de cassure]\label{def:s1:trend-following:breakout}
Une \emph{règle de cassure}\index{règle de cassure} (breakout) prend une position longue lorsque le prix d'un marché dépasse son plus haut niveau des $L$ derniers jours et une position courte
lorsqu'il passe sous son plus bas, et conserve la position jusqu'à la cassure inverse~; le canal formé par ses plus hauts et ses plus bas est aussi appelé
canal de Donchian.
\end{definition}

Les trois familles mesurent la même chose, le signe d'une pente récente lissée, avec des pondérations différentes sur le passé~: le rendement passé pondère de manière égale les $L$ derniers jours,
un croisement pondère davantage les jours récents (une pondération en triangle), et une cassure ne réagit qu'aux nouveaux extrêmes, ce qui en fait la règle la plus lente à changer de position.
Le \Cref{lst:s1:trend-following:signals} les calcule, chacune à partir des données jusqu'à la clôture du jour $t$, pour une position conservée le jour $t + 1$.

\begin{proposition}[Les gains de tendance sont convexes dans le mouvement]\label{prop:s1:trend-following:convex}
Sur une fenêtre de $T$ jours, une position égale au rendement cumulé depuis le début de la fenêtre, $w_t = \sum_{s < t} r_s$, rapporte
\[
\sum_{t=1}^{T} w_t r_t = \tfrac12\Big(\sum_{t=1}^{T} r_t\Big)^2 - \tfrac12 \sum_{t=1}^{T} r_t^2 .
\]
\end{proposition}
\begin{proof}
$\big(\sum_t r_t\big)^2 = \sum_t r_t^2 + 2\sum_{s < t} r_s r_t$, et $\sum_t w_t r_t = \sum_{s < t} r_s r_t$.
\end{proof}

Le suiveur de tendance est long du carré du mouvement de la fenêtre et court de sa variance réalisée~: il gagne lorsque les mouvements sont amples, dans un sens comme dans l'autre, et perd dans
les marchés heurtés où le cours erre sans aller nulle part. C'est la forme d'un straddle long, et c'est pourquoi les rendements de la stratégie face au
marché boursier dessinent un smile (\cref{fig:s1:trend-following:smile}). Les règles de signe sont des versions non linéaires du même pari. Les gains espérés proviennent de deux
sources~: le carré de toute dérive persistante et l'autocorrélation positive des rendements. Dans l'univers synthétique, les deux proviennent d'une dérive introduite volontairement.

\section{L'univers synthétique de contrats à terme}

\texttt{firm.synthfut} (\cref{lst:s1:trend-following:synthfut}) simule trente ans de rendements quotidiens en excès sur quarante marchés, dix dans chacune des quatre
classes~: indices actions, obligations d'État, devises et matières premières, avec des volatilités annuelles réalisées de 17.3~\%, 6.9~\%, 9.4~\% et 27.1~\%. Chaque
marché a une dérive qui suit un processus AR(1) de demi-vie d'un an et d'écart type égal à la moitié de sa volatilité~: c'est la tendance, introduite volontairement. Chacun
a aussi un carry, dont la moitié est perçue comme rendement espéré (le chapitre~20 le négocie). La volatilité se regroupe au sein de chaque classe~; les chocs suivent une loi de Student-$t$ avec un
facteur commun à la classe. Trois krachs boursiers, débutant aux années 7.5, 18 et 27 et durant chacun cent jours, ajoutent sur l'épisode une dérive régulière de $-30~\%$ aux actions, $-10~\%$
aux matières premières et $+3~\%$ aux obligations, font chuter les devises à carry élevé et augmentent la volatilité de moitié.

L'univers est conçu pour l'ensemble des chapitres consacrés aux contrats à terme~: sa vérité (dérives, carry, volatilités, dates des krachs) est conservée, et \texttt{curve} et
\texttt{contracts} donnent le logarithme du prix à terme pour toute échéance de livraison, de sorte que le chapitre~21 puisse négocier des calendar spreads sur ses courbes et que le chapitre~23 puisse repérer ses
matières premières saisonnières.

\section{Un portefeuille de tendances}

La position sur chaque marché est son signal multiplié par un notionnel visant une volatilité annuelle de 40~\% divisée par le nombre de marchés, à partir d'une prévision de volatilité
pondérée exponentiellement avec un centre de masse de soixante jours (\cref{lst:s1:trend-following:book})~; les coûts sont de 2, 1, 2 et 4 points de base par unité de
notionnel échangé pour les actions, les obligations, les devises et les matières premières (hypothèse). Le portefeuille est mesuré à partir de la deuxième année, lorsque chaque signal est défini.

\begin{center}
\small
\setlength{\tabcolsep}{4pt}%
\begin{tabular}{@{}llccc@{}}
\toprule
famille & vitesse & ratio de Sharpe après coûts & avant coûts & rotation (fois le brut par an)\\
\midrule
rendement passé & 1 mois & $-0.00$ & 0.32 & 49\\
 & 3 mois & 0.26 & 0.46 & 29\\
 & 12 mois & 0.66 & 0.76 & 15\\
moyennes mobiles & 5/20 jours & 0.10 & 0.32 & 33\\
 & 20/100 jours & 0.43 & 0.48 & 9\\
 & 50/250 jours & 0.72 & 0.76 & 5\\
cassure & 20 jours & 0.17 & 0.29 & 18\\
 & 55 jours & 0.34 & 0.39 & 8\\
 & 250 jours & 0.78 & 0.80 & 3\\
\bottomrule
\end{tabular}
\end{center}

Au sein de chaque famille, la règle lente l'emporte, car la dérive introduite est lente~; d'une famille à l'autre, à vitesse égale, les résultats sont proches, car les règles
mesurent la même pente. Les coûts départagent les règles rapides~: la règle à un mois fait tourner son brut 49 fois par an, et son ratio de Sharpe de 0.32 avant coûts
ne vaut plus rien après eux. Selon le lookback (\cref{fig:s1:trend-following:lookbacks}), le ratio de Sharpe monte jusqu'à six mois puis reste stable au-delà~: un lookback plus long
mesure la dérive avec moins de bruit mais réagit plus tard à ses changements, et, pour une demi-vie d'un an, les deux effets se compensent sur une large plage. Moskowitz,
Ooi et Pedersen ont retrouvé la même plage sur 58 marchés réels~: une persistance de un à douze mois, puis un retournement partiel au-delà.

\begin{omfigure}
\begin{tikzpicture}
\begin{axis}[width=10cm,height=5cm,xlabel={\small lookback de la règle de rendement passé (mois)},ylabel={\small ratio de Sharpe après coûts},
  tick label style={font=\footnotesize},xmin=0,xmax=25,ymin=-0.1,ymax=0.8,grid=major,grid style={gray!25},table/col sep=comma]
\addplot[omThm,very thick,mark=*] table[x=months,y=sr]{figdata/strategies-1/19-trend-following/lookbacks.csv};
\addplot[black!40,dashed] coordinates {(0,0) (25,0)};
\end{axis}
\end{tikzpicture}

\omcaption{\omterm{def:s1:trend-following:trend}{Momentum temporel} sur l'univers synthétique de quarante contrats à terme, trente ans~: le ratio de Sharpe après coûts de la règle de signe du rendement passé
pondérée par la volatilité, selon le lookback. Les dérives introduites ont une demi-vie d'un an. Données~: \texttt{s1\_trend.lookbacks}.}\label{fig:s1:trend-following:lookbacks}
\end{omfigure}

Un praticien ignore la vitesse des tendances à venir~; les portefeuilles de tendance combinent donc plusieurs vitesses. Le mélange des trois règles de rendement passé, en moyennant leurs
positions, rapporte 0.40~; celui des neuf règles, 0.55. Les deux sont inférieurs à la meilleure règle isolée, ce que coûte la vision rétrospective~: la meilleure règle n'est connue
qu'après coup, et la choisir parmi neuf est précisément la sélection que pénalisent les mesures de surapprentissage du livre~7, chapitre~20. Le mélange est une assurance contre
le fait que la vitesse de la prochaine tendance diffère de la précédente.

\section{Pondération par la volatilité}

\begin{definition}[Ciblage de volatilité]\label{def:s1:trend-following:voltarget}
Le \emph{ciblage de volatilité}\index{ciblage de volatilité} dimensionne une position, ou un portefeuille entier, de sorte que sa volatilité prévue soit égale à une cible fixe~:
le notionnel est égal à la cible divisée par la prévision, recalculée à mesure que la prévision évolue, en général avec un cap sur l'effet de levier.
\end{definition}

Les portefeuilles de tendance l'utilisent à deux niveaux. Par marché, il fait porter le même risque à un contrat à terme sur obligations et à un contrat à terme sur gaz naturel, de sorte que le portefeuille soit diversifié en
risque et non dominé par les marchés les plus volatils. Remplacer les positions pondérées par la volatilité des trois règles de rendement passé par des positions de notionnel égal
fait tomber le ratio de Sharpe du mélange de 0.40 à 0.31~: les matières premières dominent le portefeuille à notionnel égal, et les obligations pèsent à peine. Baltas et Kosowski
montrent que le choix de l'estimateur de volatilité compte aussi, surtout par la rotation du portefeuille.

Au niveau du portefeuille, un ciblage à 10~\% fait passer la volatilité réalisée du mélange de 7.7~\% à 10.2~\% et abaisse son ratio de Sharpe de 0.40 à 0.34~; le
plus grand drawdown, de 3.7 écarts-types annuels sans ciblage, atteint 3.8 avec ciblage. Dans ce modèle, la volatilité du portefeuille mesure surtout le nombre de marchés qui
suivent ensemble une même tendance~: quand ils divergent, la volatilité du portefeuille baisse et le ciblage la relève, justement quand son edge est faible. Le ciblage de portefeuille est une décision de risque,
non un alpha, et le chapitre~26 mesure quand il aide.

\section{Le long historique et l'alpha de crise}

\begin{definition}[Alpha de crise]\label{def:s1:trend-following:crisis}
L'\emph{alpha de crise}\index{alpha de crise} est la propension d'une stratégie à dégager des rendements positifs pendant les baisses prolongées du marché boursier, mesurée par son
rendement sur les plus fortes baisses du marché~; pour le \omterm{def:s1:trend-following:trend}{suivi de tendance}, il tient à ce que la stratégie soit vendeuse d'actions et positionnée pour la fuite vers la qualité une fois que la
baisse dure depuis plus longtemps que le lookback du signal.
\end{definition}

Les krachs synthétiques illustrent le mécanisme (la \cref{fig:s1:trend-following:paths}). Le premier, à l'année 7.5, s'est trouvé petit (les actions ont perdu 6.8~\% sur
ses cent jours, la baisse introduite étant compensée par le bruit) et le portefeuille de tendance, à une volatilité de 10~\%, a gagné 2.3~\%. Dans le deuxième, les actions ont chuté de 24.3~\% et le
portefeuille de tendance a gagné 21.1~\%~; dans le troisième, 35.1~\% et 39.0~\%. Les gains arrivent une fois les baisses engagées~: les signaux passent à la vente d'actions et à l'achat d'obligations quelques
semaines après, et la baisse régulière fait le reste. Un krach qui se produit en une journée, comme celui du 19 octobre 1987 (livre~6, chapitre~22), ne laisse pas le temps aux signaux d'agir, et
ne rapporte rien.

\begin{omfigure}
\begin{tikzpicture}
\begin{axis}[width=11cm,height=5.4cm,xlabel={\small année},ylabel={\small rendement cumulé (\%)},tick label style={font=\footnotesize},
  xmin=1,xmax=30,ymin=-20,ymax=160,grid=major,grid style={gray!25},legend columns=2,legend style={font=\scriptsize,at={(0.5,-0.3)},
  anchor=north,draw=none,fill=none,/tikz/every even column/.append style={column sep=8pt}},table/col sep=comma]
\fill[omExo!15] (axis cs:7.5,-20) rectangle (axis cs:7.897,160);
\fill[omExo!15] (axis cs:18,-20) rectangle (axis cs:18.397,160);
\fill[omExo!15] (axis cs:27,-20) rectangle (axis cs:27.397,160);
\addplot[omThm,thick] table[x=year,y=trend]{figdata/strategies-1/19-trend-following/paths.csv};
\addplot[omDef,thick] table[x=year,y=equity]{figdata/strategies-1/19-trend-following/paths.csv};
\legend{portefeuille de tendance (cible de volatilité de 10~\%),classe actions (moyenne de dix indices)}
\end{axis}
\end{tikzpicture}

\omcaption{Le portefeuille de tendance mélangé, ciblé à une volatilité de 10~\%, et la moyenne des dix contrats à terme sur indices actions de l'univers synthétique, en sommes cumulées
de rendements journaliers excédentaires~; les bandes grisées sont les trois krachs boursiers introduits. Données~: \texttt{s1\_trend.paths}.}\label{fig:s1:trend-following:paths}
\end{omfigure}

Trimestre par trimestre, les rendements du portefeuille de tendance face à ceux de la classe actions dessinent le «~smile~» de la proposition sur la convexité
(la \cref{fig:s1:trend-following:smile})~: une quadratique ajustée aux 115 trimestres a une courbure de 2.63, et le pire trimestre de la classe actions ($-36.2\%$) a été le
meilleur du portefeuille ($+31.6~\%$). Sur l'ensemble des trimestres, la corrélation vaut $-0.09$~: le portefeuille de tendance n'est pas une position vendeuse sur les actions, mais une position acheteuse sur les grands mouvements.

\begin{omfigure}
\begin{tikzpicture}
\begin{axis}[width=10cm,height=5.6cm,xlabel={\small rendement de la classe actions au trimestre (\%)},ylabel={\small rendement du portefeuille de tendance (\%)},
  tick label style={font=\footnotesize},xmin=-40,xmax=30,ymin=-20,ymax=40,grid=major,grid style={gray!25},legend pos=north east,
  legend style={font=\scriptsize},table/col sep=comma]
\addplot[only marks,mark=*,mark size=1.2pt,omDef] table[x=equity,y=trend]{figdata/strategies-1/19-trend-following/smile.csv};
\addplot[omThm,very thick] table[x=equity,y=trend]{figdata/strategies-1/19-trend-following/smile_fit.csv};
\legend{trimestres,ajustement quadratique}
\end{axis}
\end{tikzpicture}

\omcaption{Le smile de tendance~: le rendement trimestriel du portefeuille de tendance mélangé (cible de volatilité de 10~\%) face à celui de la classe actions, sur 115 trimestres de l'univers synthétique,
avec une quadratique ajustée (courbure 2.63). Données~: \texttt{s1\_trend.smile}.}\label{fig:s1:trend-following:smile}
\end{omfigure}

L'historique public est plus long que n'importe quel historique synthétique. Moskowitz, Ooi et Pedersen ont trouvé un portefeuille diversifié de \omterm{def:s1:trend-following:trend}{momentum temporel}, avec des
rendements anormaux substantiels, une faible exposition aux facteurs standard et sa meilleure performance sur les marchés extrêmes~; Hurst, Ooi et Pedersen ont étendu les résultats jusqu'en 1880,
avec des rendements moyens positifs sur chaque décennie et de bonnes performances dans huit des dix plus grandes périodes de crise.

Un marché peut être testé sur données réelles ici~: le NYMEX WTI crude oil, à partir des cours de règlement quotidiens de l'EIA portant sur les deux premiers contrats à terme (livre~3, chapitre~2),
reportés cinq jours avant l'échéance en une série continue. La règle des douze mois, à une cible de volatilité de 40~\%, a rapporté 9.0~\% par an de 1986 à 2024, avec
un ratio de Sharpe de 0.22, en position longue sur 54~\% des jours~; la simple détention du contrat à terme en position longue a donné un ratio de Sharpe de 0.04. Un seul marché constitue un petit échantillon (un ratio de Sharpe mesuré
sur 38 ans a une erreur type d'environ 0.16), et l'argument en faveur du \omterm{def:s1:trend-following:trend}{suivi de tendance} repose sur les nombreux marchés qui le diversifient.

\section{La décennie récente}

Un long historique ne garantit pas la décennie suivante. Baltas et Kosowski étudient la sous-performance du \omterm{def:s1:trend-following:trend}{momentum temporel} après la crise financière
mondiale et la relient au niveau des corrélations deux à deux entre marchés~: lorsque chaque marché évolue avec tous les autres, le portefeuille ne prend en réalité qu'une seule position, répétée quarante fois.
L'univers synthétique apporte l'autre moitié de la réponse, le bruit d'échantillonnage. Le portefeuille mélangé, dont l'edge introduit reste le même sur toute la période, affiche
des ratios de Sharpe de 0.56, 0.24 et 0.34 sur ses trois décennies. Le ratio de Sharpe d'une décennie a une erreur type d'environ $1/\sqrt{10} = 0.32$, si bien qu'une stratégie
avec un ratio de Sharpe véritable de 0.4 peut afficher n'importe quelle valeur entre 0.08 et 0.72 pendant la plupart des décennies. Une mauvaise décennie est un faible indice que l'edge a disparu~; un
mécanisme qui explique pourquoi il aurait disparu (corrélations, encombrement, coûts) constitue un argument plus solide.

\section{Fiches de stratégie}

\begin{strategyfile}{Momentum temporel}\label{strat:s1:trend-following:tsmom}
\sfield{Qui paie le gain, et pourquoi} Les hedgers qui négocient contre les tendances, et les investisseurs qui sous-réagissent aux nouvelles puis sur-réagissent ensuite~; les spéculateurs profitent aux
dépens des hedgers, selon les données de Moskowitz, Ooi et Pedersen.
\sfield{Instruments et plateformes} Contrats à terme liquides sur indices actions, obligations, devises et matières premières.
\sfield{Signal} Le signe du rendement excédentaire passé sur 12 mois de chaque marché.
\sfield{Dimensionnement et exécution} Positions pondérées par la volatilité, rééquilibrées chaque jour ou chaque semaine~; contrats à terme reportés avant l'échéance.
\sfield{Coûts} Faibles~: une rotation d'environ quinze fois le notionnel brut par an pour un lookback de douze mois.
\sfield{Comment elle s'éteint} Marchés corrélés~; marchés heurtés sans mouvements persistants~; encombrement.
\sfield{Horizon, capacité, infrastructure} Quelques mois~; grande capacité sur les contrats à terme liquides~; une gestion du roll et des appels de marge.
\sfield{Backtest honnête} Contrats reportés avec leurs coûts~; prévisions de volatilité connues à l'époque~; lookback non choisi a posteriori.
\sfield{Sources} Moskowitz, Ooi et Pedersen (2012)~; ce chapitre~: 0.66 après coûts sur l'univers synthétique.
\end{strategyfile}

\begin{strategyfile}{Portefeuille à croisement de moyennes mobiles}\label{strat:s1:trend-following:ma}
\sfield{Qui paie le gain, et pourquoi} Comme pour le \omterm{def:s1:trend-following:trend}{momentum temporel}.
\sfield{Instruments et plateformes} Comme pour le \omterm{def:s1:trend-following:trend}{momentum temporel}.
\sfield{Signal} Le signe de la moyenne mobile courte moins la moyenne mobile longue du cours de chaque marché, par exemple 50/250 jours.
\sfield{Dimensionnement et exécution} Positions pondérées par la volatilité~; le croisement opère moins souvent que la règle de rendement passé, à vitesse comparable.
\sfield{Coûts} Les plus faibles des trois familles aux vitesses lentes (cinq fois le notionnel brut par an pour 50/250 ici).
\sfield{Comment elle s'éteint} Comme pour le \omterm{def:s1:trend-following:trend}{momentum temporel}.
\sfield{Horizon, capacité, infrastructure} Quelques mois~; comme pour le \omterm{def:s1:trend-following:trend}{momentum temporel}.
\sfield{Backtest honnête} Le couple de fenêtres fixé avant le test, ou choisi avec les contrôles de surapprentissage du livre~7.
\sfield{Sources} Baltas et Kosowski (2020) sur les règles de négociation et la rotation~; ce chapitre~: 0.72.
\end{strategyfile}

\begin{strategyfile}{Système de cassure}\label{strat:s1:trend-following:breakout}
\sfield{Qui paie le gain, et pourquoi} Comme pour le \omterm{def:s1:trend-following:trend}{momentum temporel}~; les cassures rencontrent aussi les ordres stop placés aux plus hauts et aux plus bas récents.
\sfield{Instruments et plateformes} Contrats à terme liquides.
\sfield{Signal} Nouveaux plus hauts et nouveaux plus bas de $L$ jours (20, 55 ou 250 jours).
\sfield{Dimensionnement et exécution} Pondérés par la volatilité~; ordres placés au niveau du canal, souvent sous forme d'ordres stop.
\sfield{Coûts} Glissement aux cassures, lorsque de nombreux stops se déclenchent ensemble.
\sfield{Comment elle s'éteint} Fausses cassures dans les marchés en range~; encombrement des stops (crowding).
\sfield{Horizon, capacité, infrastructure} De quelques semaines à quelques mois.
\sfield{Backtest honnête} Exécutions au niveau de la cassure, augmentées du glissement, et non au cours de clôture de la journée.
\sfield{Sources} Aucun chiffre de performance vérifié~; ce chapitre~: 0.78 à 250 jours, 0.17 à 20 jours.
\end{strategyfile}

\begin{strategyfile}{Mélange de tendance multi-vitesses}\label{strat:s1:trend-following:blend}
\sfield{Qui paie le gain, et pourquoi} Comme pour le \omterm{def:s1:trend-following:trend}{momentum temporel}, à chaque horizon.
\sfield{Instruments et plateformes} Contrats à terme liquides.
\sfield{Signal} La moyenne des règles à plusieurs vitesses et dans plusieurs familles.
\sfield{Dimensionnement et exécution} Positions moyennées avant la négociation, de sorte que les règles opposées se compensent.
\sfield{Coûts} Ceux des composantes rapides, réduits par la compensation.
\sfield{Comment elle s'éteint} Comme pour le \omterm{def:s1:trend-following:trend}{momentum temporel}.
\sfield{Horizon, capacité, infrastructure} De quelques semaines à un an.
\sfield{Backtest honnête} L'ensemble des vitesses est fixé à l'avance~; la comparaison avec la meilleure règle isolée ne vaut qu'à titre rétrospectif.
\sfield{Sources} Hurst, Ooi et Pedersen (2017)~; ce chapitre~: 0.55 pour neuf règles, 0.40 pour trois vitesses d'une même règle.
\end{strategyfile}

\begin{strategyfile}{Portefeuille de tendance pondéré par la volatilité}\label{strat:s1:trend-following:volscaled}
\sfield{Qui paie le gain, et pourquoi} Comme pour le \omterm{def:s1:trend-following:trend}{momentum temporel}.
\sfield{Instruments et plateformes} Contrats à terme liquides sur les quatre classes d'actifs.
\sfield{Signal} Toute règle de tendance.
\sfield{Dimensionnement et exécution} Chaque marché à la même volatilité prévue~; le portefeuille, optionnellement, à une cible de volatilité avec un plafond de levier.
\sfield{Coûts} Rotation supplémentaire liée aux variations de volatilité~; des estimateurs efficaces la réduisent.
\sfield{Comment elle s'éteint} Prévisions de volatilité qui retardent sur les changements de régime.
\sfield{Horizon, capacité, infrastructure} Comme pour la règle~; une chaîne de calcul de volatilité quotidienne.
\sfield{Backtest honnête} Prévisions fondées sur les seules données passées~; plafonds de levier tels qu'ils se seraient appliqués.
\sfield{Sources} Moskowitz, Ooi et Pedersen (2012)~; Baltas et Kosowski (2020)~; ce chapitre~: 0.40 pondéré, contre 0.31 à notionnel égal.
\end{strategyfile}

\begin{strategyfile}{Tendance sur les contrats à terme WTI}\label{strat:s1:trend-following:wti}
\sfield{Qui paie le gain, et pourquoi} Les producteurs et les consommateurs qui se couvrent, et l'ajustement lent des cours du pétrole aux chocs d'offre.
\sfield{Instruments et plateformes} Contrats à terme sur pétrole brut WTI du NYMEX, roulés avant l'échéance.
\sfield{Signal} Le signe du rendement des 12 derniers mois de la série roulée.
\sfield{Dimensionnement et exécution} Une cible de volatilité de 40~\%, avec roll cinq jours avant l'échéance.
\sfield{Coûts} Le spread du roll et la forme de la courbe (le contango pénalise une position longue).
\sfield{Comment elle s'éteint} Comme il s'agit d'un seul marché, elle n'est pas diversifiée~: elle s'éteint avec le régime du pétrole.
\sfield{Horizon, capacité, infrastructure} Quelques mois.
\sfield{Backtest honnête} La série roulée, et non le cours du contrat le plus proche~; le règlement-livraison négatif d'avril 2020 traité en roulant plus tôt.
\sfield{Sources} Règlements-livraisons de l'EIA~; ce chapitre~: un ratio de Sharpe de 0.22, 1986--2024.
\end{strategyfile}

\section{Tutoriel~: alpha de crise}\label{tut:s1:trend-following:tut}

\begin{tutorial}
\textbf{Objectif.} Construire l'univers synthétique de contrats à terme, faire tourner les règles de tendance à trois vitesses dans trois familles, les mélanger et les pondérer, puis mesurer le portefeuille dans
les krachs introduits à dessein et sur le pétrole brut réel. \textbf{État final~:} le tableau et les trois figures.
\begin{tutsteps}
\item \textbf{L'univers}~: dérives, carry, états de volatilité, facteurs de classe d'actifs et krachs, jour après jour.
\omcode{code/firm/synthfut/firm_synthfut.py}{89}{103}{Une journée de l'univers synthétique de contrats à terme.}\label{lst:s1:trend-following:synthfut}
\item \textbf{Les signaux}~: prévision de volatilité, rendement passé, croisement de moyennes mobiles et cassure.
\omcode{code/firm/trendfollow/firm_trendfollow.py}{24}{60}{Signaux de tendance et prévision de volatilité.}\label{lst:s1:trend-following:signals}
\item \textbf{Le portefeuille}~: positions, P\&L net des coûts, et cible de volatilité du portefeuille.
\omcode{code/firm/trendfollow/firm_trendfollow.py}{63}{84}{Positions, P\&L et ciblage de volatilité.}\label{lst:s1:trend-following:book}
\item \textbf{Lancer} \texttt{table()}, \texttt{lookbacks()}, \texttt{summary()}, \texttt{crises()}, \texttt{smile()}, \texttt{decades()} et
\texttt{wti()}, et \texttt{fig\_trend.py}.
\end{tutsteps}
\textbf{Pour aller plus loin.} Raccourcir la demi-vie des dérives et trouver la meilleure fenêtre~; faire survenir un krach en cinq jours au lieu de cent~; remplacer
le signe par un signal continu (le rendement passé divisé par sa volatilité) et comparer la rotation.
\end{tutorial}

\section{Construction~: univers de contrats à terme et bibliothèque de tendance}\label{bld:s1:trend-following:trend}

\begin{build}
\bfield{Rôle} \texttt{firm.synthfut}~: un univers de contrats à terme déterministe avec dérives introduites, carry, regroupement de volatilité, krachs et courbes saisonnières,
pour les chapitres 19 à 25. \texttt{firm.trendfollow}~: signaux de tendance, positions pondérées par la volatilité, P\&L net des coûts et ciblage du portefeuille.
\bfield{Interface} \texttt{FutConfig(\dots)}, \texttt{simulate\_futures(cfg)}, \texttt{season(amp, phase, t)}, \texttt{curve(F, t, taus)},
\texttt{contracts(F, t, n, every)}~; \texttt{ewma\_vol}, \texttt{tsmom}, \texttt{ma\_cross}, \texttt{breakout}, \texttt{positions}, \texttt{run},
\texttt{vol\_target}.
\bfield{Règles} Les signaux utilisent les données jusqu'à la clôture de $t$ pour les positions détenues sur $t + 1$~; l'univers conserve ses valeurs vraies~; les rendements des contrats à terme ne contiennent aucune saisonnalité
(la courbe l'anticipe).
\bfield{Tests d'acceptation} \texttt{code/firm/synthfut/tests/} et \texttt{code/firm/trendfollow/tests/}~: formes et déterminisme, volatilités par classe,
la chute lors d'un krach, la pente de la courbe égale au carry~; signaux sur un chemin construit à la main, P\&L et coûts calculés à la main, la prévision de volatilité et sa cible.
\bfield{Pour aller plus loin} Signaux continus~; une famille de croisements EWMA~; calendriers de roll et P\&L au niveau du contrat.
\end{build}

\begin{omsources}
\item T.~J.~Moskowitz, Y.~H.~Ooi et L.~H.~Pedersen, «~Time series momentum~», \emph{Journal of Financial Economics} 104(2), 2012.
\item B.~Hurst, Y.~H.~Ooi et L.~H.~Pedersen, «~A century of evidence on trend-following investing~», \emph{Journal of Portfolio Management} 44(1), 2017.
\item N.~Baltas et R.~Kosowski, «~Demystifying time-series momentum strategies: volatility estimators, trading rules and pairwise correlations~», dans
\emph{Market Momentum}, Wiley, 2020.
\item US Energy Information Administration, règlements-livraisons quotidiens des contrats à terme sur pétrole brut WTI du NYMEX, contrats 1 et 2.
\end{omsources}

\section{Exercices}

\begin{exercise}[$\star$]\label{exo:s1:trend-following:1}
Une variance à pondération exponentielle de centre de masse 60 jours a un facteur de décroissance $\lambda = 60/61$. Déterminer ce facteur et
la demi-vie des poids, en jours.
\end{exercise}

\begin{exercise}[$\star$]\label{exo:s1:trend-following:2}
Un portefeuille de 40 marchés vise une volatilité de 40~\% divisée par 40 pour chaque marché. Déterminer le notionnel, en part du capital, qu'il détient
dans une matière première à 25~\% de volatilité, puis dans un contrat à terme obligataire à 6~\%.
\end{exercise}

\begin{exercise}[$\star$]\label{exo:s1:trend-following:3}
Utiliser la proposition sur la convexité pour expliquer pourquoi un suiveur de tendance perd sur un marché qui monte de 5~\% puis baisse de 5~\% à répétition.
\end{exercise}

\begin{exercise}[$\star\star$]\label{exo:s1:trend-following:4}
Le trinôme ajusté au smile est $2.63x^2 + 0.03x - 0.01$ sur les rendements trimestriels. Calculer ce qu'il prédit pour le portefeuille de
tendance sur un trimestre où les actions chutent de 30~\%.
\end{exercise}

\begin{exercise}[$\star\star$]\label{exo:s1:trend-following:5}
Pourquoi la règle à un mois, avec un ratio de Sharpe de 0.32 avant coûts, ne rapporte-t-elle rien après ceux-ci, alors que la cassure à 250 jours ne perd presque rien en coûts~?
\end{exercise}

\begin{exercise}[$\star\star$]\label{exo:s1:trend-following:6}
Un livre de \omterm{def:s1:trend-following:trend}{suivi de tendance} a un vrai ratio de Sharpe de 0.4. Quelle plage approximative de ratios de Sharpe sur une décennie faut-il attendre, et que conclure d'une décennie à 0.1~?
\end{exercise}

\begin{exercise}[$\star\star\star$]\label{exo:s1:trend-following:7}
\emph{Programmation.} Simuler l'univers avec \texttt{FutConfig(trend\_half\_life=63.0)} et relancer \texttt{lookbacks((3, 6, 12))}. Quelle fenêtre de référence
est la meilleure, et pourquoi n'est-ce pas trois mois~?
\end{exercise}

\begin{exercise}[$\star\star\star$]\label{exo:s1:trend-following:8}
\emph{Trouver la faille.} «~Nous avons testé 200 combinaisons de fenêtres de référence et de fenêtres de moyennes mobiles sur nos 40 contrats à terme, et la meilleure a produit un ratio de Sharpe de 1.3
après coûts~; nous allons la mettre en production.~»
\end{exercise}

\section{Problème~: Alpha de crise}

\begin{problem}[{Problème du week-end --- un livre de suivi de tendance à travers trois krachs}]\label{pb:s1:trend-following:1}
L'univers synthétique de contrats à terme du chapitre, la série du WTI et le bilan public.

\medskip\noindent\textbf{Partie I --- Règles.}
\begin{enumerate}
\item Définir le \omterm{def:s1:trend-following:trend}{suivi de tendance}, le \omterm{def:s1:trend-following:trend}{momentum temporel} et une \omterm{def:s1:trend-following:breakout}{règle de cassure}.
\item Écrire les trois familles de signaux et dire comment elles pondèrent le passé.
\item Énoncer et démontrer la proposition sur la convexité.
\item Quelles sont les deux sources de profit espéré d'une règle de tendance~?
\end{enumerate}

\medskip\noindent\textbf{Partie II --- L'univers et le livre.}
\begin{enumerate}[resume]
\item Décrire les dérives, les carry, la volatilité et les krachs de \texttt{firm.synthfut}.
\item Comment les positions sont-elles dimensionnées, et que supposent les coûts~?
\item Donner les ratios de Sharpe après coûts des neuf règles.
\item Pourquoi les règles lentes gagnent-elles ici, et que coûte et que rapporte le mélange~?
\end{enumerate}

\medskip\noindent\textbf{Partie III --- Mise à l'échelle et crises.}
\begin{enumerate}[resume]
\item Définir le \omterm{def:s1:trend-following:voltarget}{ciblage de volatilité}~; comparer les livres pondérés par la volatilité et les livres à notionnel égal.
\item Que fait le ciblage de portefeuille au mélange, et pourquoi~?
\item Définir l'\omterm{def:s1:trend-following:crisis}{alpha de crise} et donner les rendements du livre lors des trois krachs.
\item Décrire le smile et sa courbure.
\end{enumerate}

\medskip\noindent\textbf{Partie IV --- Le verdict.}
\begin{enumerate}[resume]
\item Énoncer le \emph{résultat nommé}~: le rendement du portefeuille de tendance lors des krachs boursiers simulés et son ratio de Sharpe selon les vitesses.
\item Que la règle à douze mois a-t-elle rapporté sur le WTI, et dans quelle mesure un seul marché suffit-il à convaincre~?
\item Que trouvent Moskowitz, Ooi et Pedersen, puis Hurst, Ooi et Pedersen~?
\item À quoi Baltas et Kosowski rattachent-ils la sous-performance récente~?
\item Quel est le niveau de bruit du ratio de Sharpe sur une décennie~?
\item Contre quel krach le \omterm{def:s1:trend-following:trend}{suivi de tendance} ne protégerait-il pas~?
\item En quoi le momentum de ce chapitre diffère-t-il de celui du chapitre~5~?
\item En une phrase~: que possède un suiveur de tendance~?
\end{enumerate}
\end{problem}

\section{Questions d'entretien}

\begin{interviewq}[$\star$ \iqroles{chercheur}]\label{iq:s1:trend-following:1}
Quelle est la différence entre le \omterm{def:s1:trend-following:trend}{momentum temporel} et le \omterm{def:s1:momentum:cross}{momentum transversal}~?
\end{interviewq}

\begin{interviewq}[$\star\star$ \iqroles{chercheur}]\label{iq:s1:trend-following:2}
Pourquoi les suiveurs de tendance pondèrent-ils leurs positions par la volatilité~?
\end{interviewq}

\begin{interviewq}[$\star\star$ \iqroles{chercheur, trader}]\label{iq:s1:trend-following:3}
Votre livre de \omterm{def:s1:trend-following:trend}{suivi de tendance} perd de l'argent depuis deux ans. Comment décidez-vous de le conserver en activité~?
\end{interviewq}

\begin{interviewq}[$\star\star$ \iqroles{développeur}]\label{iq:s1:trend-following:4}
De quoi un backtester de contrats à terme a-t-il besoin que celui d'actions n'exige pas~?
\end{interviewq}

\begin{interviewq}[$\star\star$ \iqroles{risque}]\label{iq:s1:trend-following:5}
Le \omterm{def:s1:trend-following:trend}{suivi de tendance} est-il une couverture pour un portefeuille d'actions~? Quand ne l'est-il pas~?
\end{interviewq}

\begin{interviewq}[$\star\star\star$ \iqroles{chercheur}]\label{iq:s1:trend-following:6}
Les rendements ont une dérive constante $\mu$ et un bruit indépendant de variance $\sigma^2$ par jour. Une règle linéaire retient $w_t = \sum_{k=1}^{L} r_{t-k}$. Dérivez son P\&L quotidien espéré
et le ratio de Sharpe de ce P\&L pour de petits $\mu/\sigma$.
\end{interviewq}
